% problem_id: S001-theorem-keel-mori
% source_id: S001 (Stacks Project)
% source_file: stacks-more-morphisms.tex
% statement_locator: stacks-more-morphisms.tex lines 3636--3645
% proof_locator: stacks-more-morphisms.tex lines 3647--3885
% env: theorem
% status: text-extracted-verbatim (difficulty judgement pending, written by hand)
%
% The statement and proof below are the author's own text, extracted verbatim.
% Nothing is paraphrased, shortened, or supplemented.

\begin{theorem}[Keel-Mori]
\label{theorem-keel-mori}
Let $\mathcal{X}$ be an algebraic stack. Assume
$\mathcal{I}_\mathcal{X} \to \mathcal{X}$ is finite.
Then there exists a uniform categorical moduli space
$$
f : \mathcal{X} \longrightarrow M
$$
and $f$ is separated, quasi-compact, and a universal homeomorphism.
\end{theorem}

\begin{proof}
We choose a set $I$\footnote{The reader who is still keeping
track of set theoretic issues should make sure $I$ is not too large.}
and for $i \in I$ a morphism of algebraic stacks
$g_i : \mathcal{X}_i \to \mathcal{X}$ as in
Lemma \ref{lemma-etale-local-finite-inertia}; we will use all
of the properties listed in this lemma without further mention.
Let
$$
f_i : \mathcal{X}_i \to M_i
$$
be as in Lemma \ref{lemma-well-nigh-affine-moduli-space}.
Consider the stacks
$$
\mathcal{X}_{ij} = \mathcal{X}_i \times_{g_i, \mathcal{X}, g_j} \mathcal{X}_j
$$
for $i, j \in I$. The projections $\mathcal{X}_{ij} \to \mathcal{X}_i$
and $\mathcal{X}_{ij} \to \mathcal{X}_j$ are separated
by Morphisms of Stacks, Lemma
\ref{stacks-morphisms-lemma-base-change-separated},
\'etale by Morphisms of Stacks, Lemma
\ref{stacks-morphisms-lemma-base-change-etale},
and induce isomorphisms on automorphism groups
(as in Morphisms of Stacks, Remark
\ref{stacks-morphisms-remark-identify-automorphism-groups}) by
Morphisms of Stacks, Lemma
\ref{stacks-morphisms-lemma-base-change-stabilizer-preserving}.
Thus we may apply Lemma \ref{lemma-etale-separated-over-well-nigh-affine}
to find a commutative diagram
$$
\xymatrix{
\mathcal{X}_i \ar[d]_{f_i} &
\mathcal{X}_{ij} \ar[d]_{f_{ij}} \ar[l] \ar[r] &
\mathcal{X}_j \ar[d]_{f_j} \\
M_i &
M_{ij} \ar[l] \ar[r] &
M_j
}
$$
with cartesian squares where $M_{ij} \to M_i$ and $M_{ij} \to M_j$
are separated \'etale morphisms of schemes; here we also use that $f_i$
is a uniform categorical quotient by
Lemma \ref{lemma-moduli-space-finite-affine}.
Claim:
$$
\coprod M_{ij} \longrightarrow \coprod M_i \times \coprod M_i
$$
is an \'etale equivalence relation.

\medskip\noindent
Proof of the claim. Set $R = \coprod M_{ij}$ and $U = \coprod M_i$.
We have already seen that $t : R \to U$ and $s : R \to U$ are \'etale.
Let us construct a morphism $c : R \times_{s, U, t} R \to R$
compatible with $\text{pr}_{13} : U \times U \times U \to U \times U$.
Namely, for $i, j, k \in I$ we consider
$$
\mathcal{X}_{ijk} =
\mathcal{X}_i \times_{g_i, \mathcal{X}, g_j} \mathcal{X}_j
\times_{g_j, \mathcal{X}, g_k} \mathcal{X}_k =
\mathcal{X}_{ij} \times_{\mathcal{X}_j} \mathcal{X}_{jk}
$$
Arguing exactly as in the previous paragraph,
we find that $M_{ijk} = M_{ij} \times_{M_j} M_{jk}$
is a categorical moduli space for $\mathcal{X}_{ijk}$.
In particular, there is a canonical morphism
$M_{ijk} = M_{ij} \times_{M_j} M_{jk} \to M_{ik}$
coming from the projection $\mathcal{X}_{ijk} \to \mathcal{X}_{ik}$.
Putting these morphisms together we obtain the morphism $c$.
In a similar fashion we construct a morphism $e : U \to R$
compatible with $\Delta : U \to U \times U$ and
$i : R \to R$ compatible with the flip $U \times U \to U \times U$.
Let $k$ be an algebraically closed field. Then
$$
\Mor(\Spec(k), \mathcal{X}_i) \to \Mor(\Spec(k), M_i) = M_i(k)
$$
is bijective on isomorphism classes and the same remains true after any
base change by a morphism $M' \to M$. This follows from our choice
of $f_i$ and Morphisms of Stacks, Lemmas
\ref{stacks-morphisms-lemma-universally-injective} and
\ref{stacks-morphisms-lemma-universally-injective-point}.
By construction of $2$-fibred products the diagram
$$
\xymatrix{
\Mor(\Spec(k), \mathcal{X}_{ij}) \ar[d] \ar[r] &
\Mor(\Spec(k), \mathcal{X}_j) \ar[d] \\
\Mor(\Spec(k), \mathcal{X}_i) \ar[r] &
\Mor(\Spec(k), \mathcal{X})
}
$$
is a fibre product of categories. By our choice of $g_i$ the
functors in this diagram induce bijections on automorphism groups.
It follows that this diagram induces a fibre product diagram
on sets of isomorphism classes! Thus we see that
$$
R(k) = U(k) \times_{|\Mor(\Spec(k), \mathcal{X})|} U(k)
$$
where $|\Mor(\Spec(k), \mathcal{X})|$ denotes the set
of isomorphism classes.
In particular, for any algebraically closed field $k$
the map on $k$-valued point is an equivalence relation.
We conclude the claim holds by
Groupoids, Lemma \ref{groupoids-lemma-etale-equivalence-relation}.

\medskip\noindent
Let $M = U/R$ be the algebraic space which is the quotient of the above
\'etale equivalence relation, see
Spaces, Theorem \ref{spaces-theorem-presentation}.
There is a canonical morphism $f : \mathcal{X} \to M$
fitting into commutative diagrams
\begin{equation}
\label{equation-fundamental-diagram}
\xymatrix{
\mathcal{X}_i \ar[r]_{g_i} \ar[d]_{f_i} & \mathcal{X} \ar[d]^f \\
M_i \ar[r] & M
}
\end{equation}
Namely, such a morphism $f$ is given by a functor
$$
f : \Mor(T, \mathcal{X}) \longrightarrow \Mor(T, M)
$$
for any scheme $T$ compatible with base change. Let $a : T \to \mathcal{X}$
be an object of the left hand side. We obtain an \'etale covering
$\{T_i \to T\}$ with $T_i = \mathcal{X}_i \times_\mathcal{X} T$
and morphisms $a_i : T_i \to \mathcal{X}_i$. Then we get
$b_i = f_i \circ a_i : T_i \to M_i$. Since
$T_i \times_T T_j = \mathcal{X}_{ij} \times_\mathcal{X} T$
we moreover get a morphism $a_{ij} : T_i \times_T T_j \to \mathcal{X}_{ij}$.
Setting $b_{ij} = f_{ij} \circ a_{ij}$ we find that
$b_i \times b_j$ factors through the monomorphism
$M_{ij} \to M_i \times M_j$. Hence the morphisms
$$
T_i \xrightarrow{b_i} M_i \to M
$$
agree on $T_i \times_T T_j$. As $M$ is a sheaf for the \'etale
topology, we see that these morphisms glue to a unique
morphism $b = f(a) : T \to M$. We omit the verification that
this construction is compatible with base change and we omit
the verification that the diagrams
(\ref{equation-fundamental-diagram}) commute.

\medskip\noindent
Claim: the diagrams (\ref{equation-fundamental-diagram}) are cartesian.
To see this we study the induced morphism
$$
h_i : \mathcal{X}_i \longrightarrow M_i \times_M \mathcal{X}
$$
This is a morphism of stacks \'etale over $\mathcal{X}$
and hence $h_i$ is \'etale (Morphisms of Stacks, Lemma
\ref{stacks-morphisms-lemma-etale-permanence}).
Since $g_i$ is separated, we see $h_i$ is separated
(use Morphisms of Stacks, Lemma
\ref{stacks-morphisms-lemma-compose-after-separated} and the fact
seen above that the diagonal of $\mathcal{X}$ is separated).
The morphism $h_i$ induces isomorphisms on automorphism groups
(Morphisms of Stacks, Remark
\ref{stacks-morphisms-remark-identify-automorphism-groups})
as this is true for $g_i$. For an algebraically closed field $k$
the diagram
$$
\xymatrix{
\Mor(\Spec(k), M_i \times_M \mathcal{X}) \ar[r] \ar[d] &
\Mor(\Spec(k), \mathcal{X}) \ar[d] \\
M_i(k) \ar[r] &
M(k)
}
$$
is a catesian diagram of categories and the top arrow
induces bijections on automorphism groups.
On the other hand, we have
$$
M(k) = U(k)/R(k) = U(k)/
U(k) \times_{|\Mor(\Spec(k), \mathcal{X})|} U(k) =
|\Mor(\Spec(k), \mathcal{X})|
$$
by what we said above. Thus the right vertical arrow in the
cartesian diagram above is a bijection on isomorphism classes.
We conclude that
$|\Mor(\Spec(k), M_i \times_M \mathcal{X})| \to M_i(k)$ is bijective.
Review: $h_i$ is a separated, \'etale, induces isomorphisms on
automorphism groups (as in Morphisms of Stacks, Remark
\ref{stacks-morphisms-remark-identify-automorphism-groups}), and
induces an equivalence on fibre categories over algebraically closed fields.
Hence it is an isomorphism by Morphisms of Stacks, Lemma
\ref{stacks-morphisms-lemma-etale-iso}.

\medskip\noindent
From the claim we get in particular the following:
we have a surjective \'etale morphism $U \to M$
such that the base change of $f$ is separated, quasi-compact,
and a universal homeomorphism. It follows that $f$ is separated,
quasi-compact, and a universal homeomorphism.
See Morphisms of Stacks, Lemma
\ref{stacks-morphisms-lemma-check-separated-covering},
\ref{stacks-morphisms-lemma-check-quasi-compact-covering}, and
\ref{stacks-morphisms-lemma-check-universal-homeomorphism-covering}

\medskip\noindent
To finish the proof we have to show that $f : \mathcal{X} \to M$
is a uniform categorical moduli space.
To prove this it suffices to show that given a flat morphism
$M' \to M$ of algebraic spaces, the base change
$$
M' \times_M \mathcal{X} \longrightarrow M'
$$
is a categorical moduli space. Thus we consider a morphism
$$
\theta : M' \times_M \mathcal{X} \longrightarrow E
$$
where $E$ is an algebraic space. For each $i$ we know that
$f_i$ is a uniform categorical moduli space. Hence we obtain
$$
\xymatrix{
M' \times_M \mathcal{X}_i \ar[d] \ar[r] &
M' \times_M \mathcal{X} \ar[d]^\theta \\
M' \times_M M_i \ar[r]^{\psi_i} &
E
}
$$
Since $\{M' \times_M M_i \to M'\}$ is an \'etale covering,
to obtain the desired morphism $\psi : M' \to E$ it suffices
to show that $\psi_i$ and $\psi_j$ agree over
$M' \times_M M_i \times_M M_j = M' \times_M M_{ij}$.
This follows easily from the fact that
$f_{ij} : \mathcal{X}_{ij} =
\mathcal{X}_i \times_\mathcal{X} \mathcal{X}_j \to M_{ij}$ is a uniform
categorical quotient; details omitted.
Then finally one shows that $\psi$ fits into the commutative diagram
$$
\xymatrix{
M' \times_M \mathcal{X} \ar[d] \ar[rd]^\theta \\
M' \ar[r]^\psi &
E
}
$$
because ``$\{M' \times_M \mathcal{X}_i \to M' \times_M \mathcal{X}\}$
is an \'etale covering'' and the morphisms $\psi_i$ fit into the
corresponding commutative diagrams by construction.
This finishes the proof of the Keel-Mori theorem.
\end{proof}
